%! Direct, Inverse \& Joint Variation — Unit 5, Lesson 5.7 — Group Activity
\documentclass[10pt]{article}
\usepackage{algebra2-article}
\usepackage{algebra2-boxes}
\newcommand{\TallMath}[1]{$\displaystyle #1\rule[-1.4em]{0pt}{3.2em}$}

\begin{document}

\pageheader{Unit 5, Lesson 5.7}{Group Activity}
\namepartnerperiod

\vspace{0.08in}

\begin{headlinebox}{forestbg}
\small\textbf{Constant Quotient, Constant Product.} One question decides every table:
\textbf{which row of arithmetic stays the same?} Quotients constant $\Rightarrow$ \emph{direct},
$y=kx$. Products constant $\Rightarrow$ \emph{inverse}, $y=\frac kx$. Neither constant
$\Rightarrow$ \emph{neither} --- a real answer, and a common one. Test \emph{every} row, not the
first one.
\end{headlinebox}

\vspace{0.06in}

\begin{tcolorbox}[colback=white, colframe=black!40, title=\textbf{Tier R --- Remediate},
    fonttitle=\bfseries, arc=1.5mm, left=3mm, right=3mm, top=2mm, bottom=2mm, breakable]
\textbf{Part 1 --- Run both tests.} Fill in both arithmetic rows, then circle your verdict and give
$k$ if there is one.
\begin{center}
\begin{minipage}{0.47\linewidth}\centering
\renewcommand{\arraystretch}{1.25}
\begin{tabular}{|l|c|c|c|c|}
\hline
\rowcolor{forestbg}
$x$ & $3$ & $5$ & $8$ & $10$ \\ \hline
$y$ & $12$ & $20$ & $32$ & $40$ \\ \hline
$y\div x$ & & & & \\ \hline
$x\cdot y$ & & & & \\ \hline
\end{tabular}\\[2pt]
{\small (a) \; direct \; / \; inverse \; / \; neither \quad $k=\blank{1.2cm}$}
\end{minipage}\hfill
\begin{minipage}{0.47\linewidth}\centering
\renewcommand{\arraystretch}{1.25}
\begin{tabular}{|l|c|c|c|c|}
\hline
\rowcolor{forestbg}
$x$ & $2$ & $4$ & $8$ & $16$ \\ \hline
$y$ & $24$ & $12$ & $6$ & $3$ \\ \hline
$y\div x$ & & & & \\ \hline
$x\cdot y$ & & & & \\ \hline
\end{tabular}\\[2pt]
{\small (b) \; direct \; / \; inverse \; / \; neither \quad $k=\blank{1.2cm}$}
\end{minipage}
\end{center}
\vspace{4pt}
\begin{center}
\begin{minipage}{0.47\linewidth}\centering
\renewcommand{\arraystretch}{1.25}
\begin{tabular}{|l|c|c|c|c|}
\hline
\rowcolor{forestbg}
$x$ & $1$ & $2$ & $3$ & $4$ \\ \hline
$y$ & $20$ & $17$ & $14$ & $11$ \\ \hline
$y\div x$ & & & & \\ \hline
$x\cdot y$ & & & & \\ \hline
\end{tabular}\\[2pt]
{\small (c) \; direct \; / \; inverse \; / \; neither \quad $k=\blank{1.2cm}$}
\end{minipage}\hfill
\begin{minipage}{0.47\linewidth}\centering
\renewcommand{\arraystretch}{1.25}
\begin{tabular}{|l|c|c|c|c|}
\hline
\rowcolor{forestbg}
$x$ & $2$ & $3$ & $5$ & $7$ \\ \hline
$y$ & $14$ & $21$ & $35$ & $49$ \\ \hline
$y\div x$ & & & & \\ \hline
$x\cdot y$ & & & & \\ \hline
\end{tabular}\\[2pt]
{\small (d) \; direct \; / \; inverse \; / \; neither \quad $k=\blank{1.2cm}$}
\end{minipage}
\end{center}
\vspace{2pt}
{\small In (c) the $y$-values go \emph{down}. That is not enough to make it inverse --- check the
product row and see.}

\vspace{5pt}
\textbf{Part 2 --- Find $k$, then answer.} Two steps every time.
\begin{enumerate}[label=(\alph*), leftmargin=2.1em, itemsep=5pt]
  \item $y$ varies \textbf{directly} as $x$, and $y=15$ when $x=3$. \; $k=\blank{1.2cm}$,
        equation $y=\blank{1.6cm}$. \; Find $y$ when $x=7$: \blank{1.4cm}
  \item $y$ varies \textbf{inversely} as $x$, and $y=9$ when $x=4$. \; $k=\blank{1.2cm}$,
        equation $y=\blank{1.6cm}$. \; Find $y$ when $x=6$: \blank{1.4cm}
  \item $y$ varies \textbf{directly} as $x$, and $y=30$ when $x=12$. \; $k=\blank{1.2cm}$. \;
        Find $x$ when $y=45$: \blank{1.4cm}
  \item $y$ varies \textbf{inversely} as $x$, and $y=5$ when $x=8$. \; $k=\blank{1.2cm}$. \;
        Find $x$ when $y=2$: \blank{1.4cm}
\end{enumerate}

\vspace{5pt}
\textbf{Part 3 --- Twins.} In \emph{both} of these tables, $y$ goes up by exactly $5$ every time $x$
goes up by $2$. Exactly one of them is a direct variation.
\begin{center}
\begin{minipage}{0.47\linewidth}\centering
\textbf{Twin (a)}\\[2pt]
\renewcommand{\arraystretch}{1.25}
\begin{tabular}{|l|c|c|c|c|}
\hline
\rowcolor{forestbg}
$x$ & $2$ & $4$ & $6$ & $8$ \\ \hline
$y$ & $5$ & $10$ & $15$ & $20$ \\ \hline
$y\div x$ & & & & \\ \hline
\end{tabular}\\[3pt]
{\small Verdict \blank{2.0cm}; \; equation \blank{2.0cm}}
\end{minipage}\hfill
\begin{minipage}{0.47\linewidth}\centering
\textbf{Twin (b)}\\[2pt]
\renewcommand{\arraystretch}{1.25}
\begin{tabular}{|l|c|c|c|c|}
\hline
\rowcolor{forestbg}
$x$ & $2$ & $4$ & $6$ & $8$ \\ \hline
$y$ & $8$ & $13$ & $18$ & $23$ \\ \hline
$y\div x$ & & & & \\ \hline
\end{tabular}\\[3pt]
{\small Verdict \blank{2.0cm}; \; equation \blank{2.0cm}}
\end{minipage}
\end{center}
\vspace{-2pt}
Both tables climb at the same steady rate. What made the verdicts different? \writeline
\end{tcolorbox}

\begin{tcolorbox}[colback=white, colframe=black!40, title=\textbf{Tier A --- Approaching Proficiency},
    fonttitle=\bfseries, arc=1.5mm, left=3mm, right=3mm, top=2mm, bottom=2mm, breakable]
\textbf{Part 1 --- A spring.} Hang a weight on a spring and it stretches. Let $w$ be the weight in
pounds and $s$ the stretch in centimeters.
\begin{center}
\renewcommand{\arraystretch}{1.3}
\begin{tabular}{|l|c|c|c|c|}
\hline
\rowcolor{forestbg}
$w$ (lb) & $4$ & $6$ & $10$ & $15$ \\ \hline
$s$ (cm) & $3$ & $4.5$ & $7.5$ & $11.25$ \\ \hline
$s\div w$ & & & & \\ \hline
\end{tabular}
\end{center}
\begin{enumerate}[label=\arabic*., leftmargin=2.1em, itemsep=5pt]
  \item Verdict: \blank{2.2cm} variation. \; $k=\blank{1.2cm}$. \; Equation: $s=\blank{2.0cm}$
  \item What are the \emph{units} of $k$? \blank{3.4cm} \; Say what $k$ means in one sentence.
        \blank{5.0cm}
  \item A $20$-pound weight stretches the spring \blank{1.6cm} cm. \; A stretch of $9$ cm means a
        weight of \blank{1.6cm} pounds.
  \item The graph of $s=kw$ passes through $(0,0)$. Explain why the \emph{situation} requires that.
        \writeline
\end{enumerate}

\vspace{5pt}
\textbf{Part 2 --- Every rectangle with area $36$.} A rectangle has base $b$ and height $h$, and its
area is always $36$ square inches.
\begin{center}
\begin{minipage}{0.44\linewidth}\centering
\begin{tikzpicture}[scale=0.30]
  \draw[step=1, linegray!40, very thin] (0,0) grid (13,14);
  \draw[->, semithick, charcoal] (0,0)--(13.6,0) node[below right, font=\scriptsize]{$b$};
  \draw[->, semithick, charcoal] (0,0)--(0,14.6) node[left, font=\scriptsize]{$h$};
  \node[below, font=\tiny] at (3,0) {$3$};
  \node[below, font=\tiny] at (6,0) {$6$};
  \node[below, font=\tiny] at (9,0) {$9$};
  \node[below, font=\tiny] at (12,0) {$12$};
  \node[left, font=\tiny] at (0,3) {$3$};
  \node[left, font=\tiny] at (0,6) {$6$};
  \node[left, font=\tiny] at (0,9) {$9$};
  \node[left, font=\tiny] at (0,12) {$12$};
  \begin{scope}
    \clip (0,0) rectangle (13,14);
    \draw[very thick, forest, domain=2.6:13, samples=170, smooth] plot (\x,{36/(\x)});
  \end{scope}
  \fill[charcoal] (3,12) circle (3.4pt);
  \fill[charcoal] (6,6) circle (3.4pt);
  \fill[charcoal] (12,3) circle (3.4pt);
\end{tikzpicture}
\end{minipage}\hfill
\begin{minipage}{0.52\linewidth}
\small
\begin{enumerate}[label=(\alph*), leftmargin=1.9em, itemsep=4pt]
  \item $b$ and $h$ vary \blank{1.8cm}, because $b\cdot h$ is always \blank{1.0cm}. Equation:
        $h=\blank{1.8cm}$
  \item Read the graph: when $b=6$, $h=\blank{1.0cm}$. What is special about that rectangle?
        \blank{3.0cm}
  \item Solve for the base when $h=8$: $b=\blank{1.4cm}$. Is that point on the drawn curve?
        \blank{1.0cm}
  \item The curve gets very close to the $b$-axis on the right but never touches it. Say what that
        means about rectangles. \blank{4.4cm}
  \item The equation is also happy with $b=-4$. Why does the situation say no?
        \blank{4.4cm}
\end{enumerate}
\end{minipage}
\end{center}

\vspace{5pt}
\textbf{Part 3 --- Error analysis.} A student is handed this table and writes what follows.
\begin{center}
\renewcommand{\arraystretch}{1.25}
\begin{tabular}{|l|c|c|c|c|}
\hline
\rowcolor{forestbg}
$x$ & $2$ & $4$ & $6$ & $8$ \\ \hline
$y$ & $6$ & $10$ & $14$ & $18$ \\ \hline
\end{tabular}
\end{center}
\vspace{2pt}
\begin{center}
\fbox{\parbox{0.80\linewidth}{\centering \vspace{2pt}
``$x=2$ goes with $y=6$, so $k=\frac{6}{2}=3$. \\
The equation is $y=3x$. \\
It's a direct variation with $k=3$.''\vspace{2pt}}}
\end{center}
\begin{enumerate}[label=(\alph*), leftmargin=2.1em, itemsep=5pt]
  \item Is the student's arithmetic on the \emph{first} pair correct? \blank{1.0cm}
  \item Test their equation on the rest of the table: $y=3x$ predicts \blank{1.0cm} at $x=4$, but
        the table says \blank{1.0cm}.
  \item Finish both rows properly. \; $y\div x$: \blank{4.4cm} \\
        $x\cdot y$: \blank{4.4cm} \; So the correct verdict is \blank{2.0cm}.
  \item The table is $y=2x+2$. Write the one-sentence rule the student needed. \writelines{2}
\end{enumerate}
\end{tcolorbox}

\begin{tcolorbox}[colback=white, colframe=black!40, title=\textbf{Tier E --- Extension},
    fonttitle=\bfseries, arc=1.5mm, left=3mm, right=3mm, top=2mm, bottom=2mm, breakable]
\textbf{Part 1 --- Two variables at once} \emph{(joint and combined variation --- enrichment)}.
\begin{enumerate}[label=(\alph*), leftmargin=2.1em, itemsep=5pt]
  \item $y$ varies jointly as $x$ and $z$, and $y=48$ when $x=4$ and $z=2$. \; $k=\blank{1.2cm}$,
        equation $y=\blank{1.6cm}$. \; Find $y$ when $x=5$, $z=3$: \blank{1.4cm}
  \item The volume $V$ of a cylinder varies jointly as its height $h$ and the \emph{square} of its
        radius $r$. When $r=5$ and $h=4$, $V=100\pi$. \\
        Then $k=\blank{1.2cm}$ and the equation is $V=\blank{2.0cm}$. Do you recognize it?
        \blank{2.6cm}
  \item $y$ varies directly as $x$ and inversely as $z$, and $y=14$ when $x=7$ and $z=3$. \;
        $k=\blank{1.2cm}$, equation $y=\blank{1.8cm}$. \; Find $y$ when $x=10$, $z=4$:
        \blank{1.4cm}
\end{enumerate}

\vspace{5pt}
\textbf{Part 2 --- No numbers required.} Fill the table using the equations alone.
\begin{center}
\renewcommand{\arraystretch}{1.5}
\begin{tabularx}{\linewidth}{@{}l l X@{}}
\hline
\textbf{Relationship} & \textbf{What you do to $x$ (and $z$)} & \textbf{What happens to $y$} \\
\hline
$y=kx$ & triple $x$ & \blank{5.0cm} \\
$y=kx$ & cut $x$ in half & \blank{5.0cm} \\
$y=\dfrac{k}{x}$ & triple $x$ & \blank{5.0cm} \\
$y=\dfrac{k}{x}$ & cut $x$ in half & \blank{5.0cm} \\
$y=kxz$ & double $x$ \emph{and} triple $z$ & \blank{5.0cm} \\
$y=\dfrac{kx}{z}$ & double $x$ \emph{and} double $z$ & \blank{5.0cm} \\
\hline
\end{tabularx}
\end{center}
\vspace{-2pt}
Prove one of them instead of trusting it. Start from $y=\dfrac{k}{x}$ and replace $x$ with $3x$:
\[ \text{new } y=\frac{k}{\rule{1.2cm}{0.4pt}}=\frac{1}{\rule{0.6cm}{0.4pt}}\cdot
   \frac{k}{x}=\frac{1}{\rule{0.6cm}{0.4pt}}\cdot(\text{old } y) \]

\vspace{5pt}
\textbf{Part 3 --- A direct variation in disguise, and a design task.}
\begin{enumerate}[label=(\alph*), leftmargin=2.1em, itemsep=5pt]
  \item Ravi says: ``$y=\dfrac{6}{x}$ is an \emph{inverse} variation between $y$ and $x$ --- but it is
        a \emph{direct} variation between $y$ and $\frac1x$.'' Test him.
        \begin{center}
        \renewcommand{\arraystretch}{1.3}
        \begin{tabular}{|l|c|c|c|c|}
        \hline
        \rowcolor{forestbg}
        $x$ & $1$ & $2$ & $3$ & $6$ \\ \hline
        $\frac1x$ & $1$ & $0.5$ & & \\ \hline
        $y=\frac6x$ & $6$ & $3$ & & \\ \hline
        $y\div\frac1x$ & & & & \\ \hline
        \end{tabular}
        \end{center}
        \vspace{-2pt}
        Is Ravi right? \blank{1.2cm} \; What is the constant? \blank{1.2cm} \; So the graph of $y$
        against $\frac1x$ would be a \blank{3.0cm}.
  \item Write an \emph{inverse} variation whose graph passes through $(4,5)$: \blank{3.0cm} \\
        Write a \emph{direct} variation whose graph passes through $(4,5)$: \blank{3.0cm}
  \item Now try to write a direct variation whose graph passes through $(0,3)$. \blank{3.0cm} \\
        Explain what goes wrong, in one sentence. \writeline
  \item Every inverse variation $y=\frac kx$ is the Lesson 5.4 parent $\frac1x$ after one
        transformation. Which one? \blank{4.6cm} \; And its asymptotes are always
        $x=\blank{0.8cm}$ and $y=\blank{0.8cm}$, no matter what $k$ is --- explain why $k$ cannot
        move them. \writeline
\end{enumerate}
\end{tcolorbox}

\end{document}
